\scanpage{91}
\mspage{90}

\diagram{\begin{tikzcd}[ampersand replacement=\&,row sep=small,column sep=small]
Y'\arrow[r,phantom,"\ll"{description}]\&Y\arrow[r,phantom,"\ll"{description}]\arrow[d,phantom,"\ll"{description,rotate=90}]\&X\arrow[d,phantom,"\ll"{description,rotate=90}]\\
\&Y_{\partial X}\arrow[r,phantom,"\ll"{description}]\arrow[d,phantom,"\ll"{description,rotate=90}]\&\partial X\\
\&\partial Y\&
\end{tikzcd}}

Or
\[
\begin{aligned}
&\deleted{\unsure{Y'_{\partial X}}}Y'|\partial X\quad(\text{induction dans }X)=\\
&\deleted{Y'_{\partial X}}Y'|(Y_{\partial X}\deleted{\unsure{\vert\partial X}})\quad(\text{induction dans }Y)
\end{aligned}
\]
\deleted{\unsure{grâce à}} \deleted{\unsure{l'hypothèse} :}

Comme $Y_{\partial X}\leq\partial Y$ (par hyp. $Y\overset{\circ}{\ll}X$), on a $Y'|Y_{\partial X}\leq Y'|\partial Y$, \deleted{\unsure{par}} or $Y'|\partial Y\unsure{=|\simeq}Y'_{\partial Y}=\varnothing$ par hyp. $Y'\in Y^\circ$, d'où $Y'|Y_{\partial X}=\varnothing$ qed.

Supposons $Y^\circ\subset X^\circ$ prouvons $Y\overset{\circ}{\ll}X$, i.e. qu'on n'a \uline{pas} $Y\ll\partial X$. Il faut utiliser ici
\[
\boxed{Y^\circ\ne\varnothing}
\]
(ce qui résulte de l'axiome de divisibilité L1, compte tenu de L3, et est de nature analogue à L1\ldots) Soit $Z\in Y^\circ$, donc $Z\in X^\circ$ par $Y^\circ\subset X^\circ$, donc $Z\parallel\partial X$. Si on avait $Y\ll\partial X$ d'où $Z\ll\partial X$, on aurait $Z\parallel Z$ par At D$_0$, absurde.
% slanted
\margin{\uline{NB} Moyennant des hyp. de \unsure{dimension finie}, la condition ``$Y^\circ\ne\varnothing\quad\forall Y\in\unsure{M|\mathcal M}$'' \uline{\unsure{équivaut}} à L1\ldots}

b) Supposons $Y\in X^\circ$, prouvons $\forall Y'\in\widetilde Y,Y'^\circ\subset X^\circ$. En \deleted{\unsure{effet}} fait, non seulement $Y'^\circ\subset\operatorname{Omb}(Y')$ est $\subset X^\circ$, mais $\operatorname{Omb}(Y')\subset\operatorname{Omb}(Y)$ l'est, \deleted{\unsure{i.e.}} et en fait \deleted{On} On a
\[
Y\in X^\circ\Rightarrow\operatorname{Omb}(Y)\subset X^\circ
\]
i.e. $X^\circ$ est une partie de $\deleted{\unsure{\operatorname{Omb}(X)}}\inserted{M}$ fermée pour $\ll$, comme on a \unsure{noté} déjà ant. Donc $\Rightarrow$ de b) est triv. Prouvons $\Leftarrow$, i.e.
